\addtocontents{toc}{\protect\Needspace{3\baselineskip}}
\originalchapter{8.821 课程详讲：面积熵界与状态计数}
本章对应第 6 讲前半的全息原理. 后半引入矩阵图, 与第 7--9 讲合在下一章展开. 前置是上一章的黑洞面积熵、广义第二定律和量子熵定义. 讨论取 $\hbar=c=\kb=1$, 因而 $\ell_P^2=G_N$.

\originalsection{为何从黑洞推向一个区域的最大熵}
普通不含引力的格点系统中, 一份局部自由度通常贡献有限熵, 整个盒子的最大熵随体积增长. 黑洞则给出与表面面积成正比的熵. 这两个计数规律之间的张力, 比“面积比体积少一维”的表面描述更具体: 需要问在一个固定几何区域里, 哪些状态仍能被容许为不改变该区域为黑洞的物理态.

取渐近平直、近似球形的孤立系统, 它被半径 $R$ 的球包住, 边界面积 $A=4\pi R^2$, 总能量 $E$ 与熵 $S_0$. 同面积 Schwarzschild 黑洞质量为 $M_A=R/(2G_N)$. 在这类近静态估计中, 若系统已满足 $2G_NE\gtrsim R$, 原先“盒内普通系统”的描写就不能继续当作一个弱引力构型. 所以先假定 $E<M_A$.

设能向系统加入总能量 $M_A-E$, 使最终黑洞面积仍对应 $A$, 且加入物质的熵 $S'\ge0$. 若这样的过程满足广义第二定律且最终外部残留熵可以按所述方式控制, 就有
\[
S_0+S'\le S_{\rm BH}(A)=\frac{A}{4G_N},
\qquad S_0\le\frac{A}{4G_N}.
\]
这就是课程用来激发全息原理的球形熵界论证. 它采用了过程、边界和回馈假设, 并没有证明任意宇宙学切片上任意体积区域的熵都服从相同的简单空间型面积界. 宇宙膨胀、强引力构型、非球形区域及可用的光片条件, 需要比这段启发更精确的表述.

这里也不是说系统能量达到 $M_A$ 的过程一定能维持原球形边界而没有任何辐射. 这个可行过程是假设的一部分, 不能把尚需控制的演化偷偷称作普适证明. 本课程随后用 AdS/CFT 的具体自由度数量级来检验这条思路, 而不是从一段黑洞吞并论证得到所有对偶细节.

\originalsection{从熵界到有效 Hilbert 空间}
一个 $D$ 维有限 Hilbert 空间中, 最大 von Neumann 熵是 $\log D$. 证明可用相对熵的非负性:
\[
D(\rho\Vert I/D)=\operatorname{Tr}\rho(\log\rho-\log(I/D))
=-S(\rho)+\log D\ge0.
\]
等号当且仅当 $\rho=I/D$. 因而若在给定能量和引力约束下的可容许态空间有效维数为 $D_{\rm eff}$, 面积界暗示
\[
\log D_{\rm eff}\lesssim\frac A{4\ell_P^2}.
\]
有效二值自由度个数 $n_{\rm bits}$ 则由 $D=2^{n_{\rm bits}}$ 给
\[
n_{\rm bits}\lesssim\frac A{4\ell_P^2\log2}.
\]
所以“每个 Planck 面积一个自由度”是数量级说法, 不能把位数、自然对数的熵和 Hilbert 空间维数三个量都直接写成 $A/(4\ell_P^2)$.

有效条件十分重要. 一个谐振子的 Hilbert 空间无限维, 但在 $n\omega\le E_{\max}$ 的约束下只有 $1+\lfloor E_{\max}/\omega\rfloor$ 个能级. 量子场论还有空间截止与能量约束, 而引力进一步排除导致坍缩的构型. 因此这里的态计数并不意味着每个局部数学算符空间都严格只有一个有限维矩阵代数, 也不是直接数未经约束的全部场论模式.

\begin{example}
半径 $R$ 区域若容许所有 $n$ 个自旋二值态, 则最大熵为 $n\log2$. 固定空间格距 $a$, $n\sim V/a^3\sim R^3/a^3$. 与面积熵相比
\[
\frac{S_{\rm lattice}}{S_{\rm BH}}
\sim\frac{R\ell_P^2}{a^3}
\]
忽略球形常数和 $\log2$. 对足够大 $R$, 朴素格点全部状态数会超过黑洞界. 由此不能宣布实际强引力体系有一个现成的反例, 因为还没核对这些格点态的能量、相互作用与自引力是否允许把它们同时当作固定几何中的普通状态. 正是这个遗漏使面积计数成为改变容许状态的线索.
\end{example}

\originalsection{辐射气体怎样提供可计算比较}
对有限 $Z$ 种自由无质量标量, $u=(\pi^2Z/30)T^4$, $s=4u/(3T)$. 体积 $V\sim R^3$ 中消去温度, 熵尺度为
\[
S_{\rm rad}\sim Z^{1/4}V^{1/4}E^{3/4}
\sim Z^{1/4}(ER)^{3/4}.
\]
若在坍缩阈值前以 $E\sim R/G_N$ 作最乐观估计,
\[
S_{\rm rad}\sim Z^{1/4}\left(\frac{R^2}{G_N}\right)^{3/4},
\qquad S_{\rm BH}\sim\frac{R^2}{G_N}.
\]
两者比为 $Z^{1/4}(\ell_P/R)^{1/2}$, 对宏观 $R$ 与通常有限种数远小于 1. 即使普通物质尽量装入能量, 在这个模型中仍达不到黑洞的面积熵. PS2.4 给出全部球形常数、固定能量密度的最大半径, 并讨论极大种数引出的有效引力截止问题.

该比较说明黑洞不能简单当作把同样普通粒子挤得更密就得到的气体. 但若把物种数 $Z$ 视为可以无限增大而不改变引力有效理论, 同一公式会形式上突破界; 这也是必须把物种数、截止、引力耦合与半经典适用范围一起核查的原因. 面积熵界不是独立于所有理论参数的随意计数法则.

\originalsection{全息原理已经告诉了什么}
在这里, 全息原理提出的是引力体系的最大可容许信息随包围区域的边界面积计数, 因而\textbf{可能}存在更少维度的描述. 它没有单凭一个不等式告诉我们边界理论的具体 Lagrangian、场数、相关函数, 或哪个边界态对应哪个内部黑洞微观态.

下一步要从规范场论一侧找更具体的结构. 若场是 $N\times N$ 矩阵, 普通图的颜色求和能给 $N^2$ 数量级的自由度, 图又自然形成二维带状曲面. 这个 $N^2$ 与引力面积系数将来会在 $G_N$ 的字典中相遇. 它们现在还是分别建立的两条证据, 不能提前把共同的数量级当成已经证明的等价.
